\scanpage{51}
Dans le cas où $k$ est de car. 0, cela nous donne une connexion intégrable absolue sur $\underline{H}^*$ (Manin n'utilise pas ce langage géométrique).
Puis, dans le cas du $H^1$, et lorsque \deleted{$X$ est} provient d'un modèle \inserted{propre} lisse, Manin montre de façon assez compliquée que
$H^1(X,\Omega^\bullet_{X/k})$ (qui est $\subset H^1_{\mathrm{DR}}(X/k)$) est stable par les opérateurs $\Theta_X$.

3.5. En fait, par des arguments très simples directs, \deleted{qui} reposant sur la ``formule d'homotopie'' de Cartan
\[
d i_X+i_X d=\Theta_X
\]
on peut montrer, pour un morphisme $f:X\longrightarrow S$:

(i) Pour toute dérivation $\xi$ de $\underline{\Theta}_S$ qui se remonte, \deleted{\unsure{$X$}}, \inserted{loc. au dessus de $S$,} en une dérivation $\bar{\xi}$ de $\underline{\Theta}_X$ [condition tjrs satisfaite si $X$ est \deleted{\unsure{à}} \uline{formellement lisse sur $S$} et \uline{affine} sur $S$]
l'effet de $(\bar{\xi},\xi)$ opérant sur $\deleted{\unsure{\mathbb{R}}}f_*(\underline{\Omega}^\bullet_{X/S})$, est indépendant \uline{mod homotopie} \deleted{de} du choix de $\bar{\xi}$, donc dépend seulement de $\xi$; donc en particulier, \inserted{si $f$ form$^{\mathrm{t}}$ lisse et affine,} les
$\mathcal{H}^i(f_*(\Omega^\bullet_{X/S}))\unsure{=|\equiv}\mathbb{R}^if_*(\Omega^\bullet_{X/S})$
sont des Modules sur lesquels les dérivations de $S$ opèrent, --- et on vérifie les lois de linéarité et de crochet facilement.
Si $S$ lisse sur un $T$, ceci donne sur $\unsure{f_*|\mathbb{R}f_*}(\Omega^\bullet_{X/S})$ et sur les $\mathbb{R}^if_*(\Omega^\bullet_{X/S})$ des conn. intégrables rel/$T$.

\margin{50}
